\begin{Exercise}[label={ex:shift}]
    Scrivere un programma che legge una matrice 3x4 ed effettua uno shift per righe verso il basso e stampa il risultato. L’ultima riga prende la posizione della prima riga.
    Scrivere le funzioni:   
    \begin{itemize}
        \item leggiArray(): che legge e restituisce un array di 4 elementi
        \item stampaArray(): che stampa un array di 4 elementi
        \item copiaArray(): che copia un array di 4 elementi in un altro array di 4 elementi
        \item leggiMatrice(): che legge e restituisce una matrice 3x4
        \item stampaMatrice(): che stampa una matrice 3x4
        \item shiftMatrice(): che effettua lo shift della matrice verso il basso
    \end{itemize}
\end{Exercise}

\begin{Answer}[ref={ex:shift}]
  	\lstinputlisting[style=cstylesmall]{src/shift.c}
\end{Answer}